This study proposed a dynamic rock--TBM spatial relation graph for component-level geological anomaly indexing during TBM excavation. The graph represents TSP-derived rock voxels and component-labelled TBM surface nodes as spatial entities, constructs candidate rock--TBM spatial relations under geometric and excavation-state constraints, and computes component-level geological anomaly indices from the resulting edge-weighted relations.

An event-centred evaluation was conducted using a real shield-sticking event in the Boshulaling Tunnel. The results show that the graph places TSP-derived anomalies, TBM components and monitoring responses in a common excavation coordinate system; reconstructs how low-velocity rock voxels become spatially associated with shield components during event development; and traces the shield-sticking interpretation back to specific rock voxels, TBM surface nodes and candidate edges. Compared with global geological summaries and pooled TBM summaries, the component-level geological anomaly index preserves the rock--machine relation structure needed for shield-specific interpretation.

The proposed graph should be understood as a spatial data representation and event-centred interpretation tool. It does not estimate contact force, produce calibrated risk probabilities or prove universal failure prediction. Its primary value is to make the spatial relations between anomalous geological units and moving TBM components explicit, component-level and traceable.